% ============================================================
%  SECCIÓN 1 – Introducción y Motivación
% ============================================================

\section{Introducción y Motivación}

\subsection{Contexto industrial: empresa AUTOELEC}

La empresa \textbf{AUTOELEC} se dedica al asesoramiento, venta, reparación y alquiler de autoelevadores y baterías. Dentro de su proceso productivo, una tarea crítica es el \textbf{ensamble de paquetes de baterías de plomo-ácido} destinadas a autoelevadores industriales. Estos paquetes se construyen ubicando manualmente múltiples vasos (celdas individuales) dentro de un rack metálico, conformando la batería final que alimenta el vehículo.

\begin{figure}[H]
    \centering
    \begin{minipage}[t]{0.47\textwidth}
        \centering
        \includegraphics[width=\linewidth]{Imagenes/Celdas_Carga.png}
        \caption{Vasos de plomo-ácido dispuestos en rack previo al ensamble.}
        \label{fig:celdas_carga}
    \end{minipage}
    \hfill
    \begin{minipage}[t]{0.47\textwidth}
        \centering
        \includegraphics[width=\linewidth]{Imagenes/Bateria.png}
        \caption{Paquete de baterías terminado, marca Eternity -- distribuido por AUTOELEC.}
        \label{fig:bateria}
    \end{minipage}
\end{figure}

\subsection{Descripción del problema}

El proceso de ensamble de vasos cuyo peso varía entre 20 y 100\,kg requiere un control de movimientos eficiente y una capacidad operativa que minimice errores, reduzca los tiempos de producción y evite daños derivados de una manipulación inadecuada de los materiales. Actualmente, al tratarse de un proceso \textbf{completamente manual}, se generan riesgos ergonómicos que provocan un desgaste prematuro en la salud del operador, quien está expuesto a esfuerzos físicos considerables y repetitivos.

El operador trabaja con un rack destinado al montaje de múltiples vasos de plomo-ácido. Para efectuar el ensamble, debe elevar manualmente cada vaso ---asistido únicamente por su propia fuerza física--- y ubicarlo en una posición determinada dentro del rack. Cabe destacar que, una vez instaladas la mayoría de las unidades, el espacio libre en el rack resulta muy reducido, lo que dificulta las maniobras de inserción y obliga al operador a trabajar en posiciones forzadas que pueden causar accidentes.

\subsection{Solución propuesta en Robótica II: manipulador electrohidráulico}

Para abordar esta problemática, en la asignatura Robótica II se propuso el desarrollo de un \textbf{Manipulador de Asistencia Ergonómica} de tipo electrohidráulico. Este concepto soporta el peso de la carga y permite al operador guiar el movimiento con un esfuerzo mínimo, es decir, un \textit{manipulador industrial ingrávido}.

De las alternativas evaluadas (neumático, eléctrico y contrapeso), la opción seleccionada fue un \textbf{brazo articulado con contrapeso y sistema electrohidráulico} de simple efecto. En esta configuración:

\begin{itemize}
    \item El contrapeso equilibra estáticamente el extremo distal del brazo.
    \item Un motor eléctrico acciona una bomba hidráulica que eleva la carga a velocidad constante al presionar el pulsador de elevación.
    \item Al activar el pulsador de descenso, el retorno del fluido por gravedad baja la carga de forma controlada.
    \item El control está integrado en una empuñadura ergonómica en el extremo distal, permitiendo que el operador trabaje adyacente a la pieza durante todo el ciclo.
\end{itemize}

Este diseño permite \textbf{acelerar significativamente los ciclos de manipulación} en comparación con el proceso manual, reduciendo a la vez la fatiga física y el riesgo de lesiones del operador.

\subsection{Motivación para el entorno de Realidad Virtual}

La teleoperación y el aprendizaje de sistemas de manipulación robótica presentan desafíos particulares: existen \textit{riesgos operativos} durante la fase de entrenamiento (colisiones con el rack, caída de vasos, golpes al operador), se requiere \textit{precisión manual} en el posicionamiento de piezas pesadas y hay una \textit{alta variabilidad} en la destreza inicial de los distintos operadores.

Un entorno de \textbf{Realidad Virtual (VR)} permite abordar estos desafíos de forma integral:

\begin{itemize}
    \item \textbf{Seguridad}: el operador aprende a manejar el manipulador sin riesgo de dañar equipos ni lastimarse, al interactuar con una réplica virtual del sistema real.
    \item \textbf{Reducción de costos}: se eliminan los costos asociados a errores durante el entrenamiento (materiales dañados, paradas de producción).
    \item \textbf{Iteración rápida}: el entorno virtual puede reconfigurarse para distintas secuencias de ensamble, cargas o restricciones sin intervención física.
    \item \textbf{Métricas objetivas}: es posible instrumentar automáticamente indicadores de desempeño como tiempo por tarea, suavidad del movimiento (\textit{jerk}), colisiones y reintentos.
    \item \textbf{Transferencia de expertise}: los patrones de movimiento aprendidos en VR pueden transferirse al control del robot real a través del mismo joystick físico utilizado durante el entrenamiento.
\end{itemize}

En este contexto, el presente proyecto propone el desarrollo de un \textbf{entorno VR para entrenar operarios} en la teleoperación del manipulador electrohidráulico de AUTOELEC, utilizando un joystick diseñado ad hoc sobre \textit{Raspberry Pi Pico} como interfaz de control, y el motor gráfico \textit{Unity} como plataforma de simulación y experiencia inmersiva.
